\documentclass[12pt,a4paper]{article}

\usepackage[utf8]{inputenc}
\usepackage[T1]{fontenc}
\usepackage{lmodern}
\usepackage[spanish]{babel}
\usepackage[margin=2.5cm]{geometry}
\usepackage{enumitem}
\IfFileExists{microtype.sty}{\usepackage{microtype}}{}
\usepackage[hidelinks]{hyperref}

\setlength{\emergencystretch}{2em}
\setlength{\parindent}{0pt}
\setlength{\parskip}{0.6em}
\setlist[enumerate]{leftmargin=*, itemsep=0.7em}
\setlist[itemize]{leftmargin=*, itemsep=0.5em}

\begin{document}

\section{Cuentas, notificaciones y auditoría de Ahorro Futuro}

\subsection{Notas del bloque A}

\subsubsection{Interacciones del diagrama}

\begin{enumerate}
    \item \textbf{A1 y A2. Cliente móvil/web y terminal de ventanilla al Servicio de cuentas.}
    RPC síncrona mediante gRPC (gRPC-Web en el navegador) sobre HTTP/2 con TLS. Envían \texttt{ConsultarSaldo(idCuenta)} y reciben el saldo. El socio y el cajero esperan la respuesta, por lo que la comunicación es síncrona \cite{birrell1984}.

    \item \textbf{A3. Servicio de cuentas al Core heredado.}
    Solicitud síncrona por la interfaz nativa del core, a través de su adaptador, con \texttt{ObtenerSaldo} y movimientos. Lleva un plazo de 1,5 s para que el socio reciba un error claro y no una espera indefinida \cite{grpcdeadlines}. El servicio de cuentas no guarda saldo propio.

    \item \textbf{A4a y C6. Servicio de cuentas y Servicio de transferencias al exchange \texttt{operaciones}.}
    Publicación asíncrona por AMQP 0-9-1 en RabbitMQ, con tipo \textit{topic}. Cuentas publica \texttt{ConsultaRealizada}. Transferencias publica \texttt{OperaciónConfirmada} después de que el core confirma. El publicador del bloque D envía los eventos de retiro y reverso mediante D8 al mismo exchange, con los estados confirmados que se detallan en sus notas.

    \item \textbf{A5. Cola \texttt{notificaciones.q} al Servicio de notificaciones.}
    Entrega asíncrona desde una cola durable. El servicio confirma el mensaje de forma manual solo tras procesarlo \cite{rabbitmqconfirms}.

    \item \textbf{A6. Cola \texttt{auditoria.q} al Servicio de auditoría.}
    Entrega asíncrona desde una cola \textit{quorum}, replicada entre nodos del broker \cite{rabbitmqquorum}. El servicio confirma después de guardar el evento. Los mensajes que no puede procesar pasan por un \textit{dead letter exchange} a \texttt{auditoria.dlq} \cite{rabbitmqdlx}.

    \item \textbf{A7. Servicio de notificaciones al proveedor push/SMS/correo.}
    Solicitud REST síncrona sobre HTTPS, con reintentos, fuera del flujo de la operación. Envía \texttt{aviso(idOperacion, canal)}.

    \item \textbf{A8. Servicio de auditoría a la Base de auditoría.}
    Acceso a datos sobre TCP con el controlador de PostgreSQL. Inserta el evento y no actualiza ni borra registros.
\end{enumerate}

\subsubsection{Justificación del diseño}

\begin{enumerate}
    \item \textbf{R1 y R2. Consulta de saldo.}
    El saldo oficial solo existe en el core. Guardar uno en cuentas arriesgaría mostrar un saldo que no refleja un débito reciente. Se usa RPC síncrona porque alguien espera la respuesta. Se descartó la mensajería, que añade un salto, y se prefirió gRPC frente a REST con JSON por su contrato y su gestión explícita de plazos; un cliente REST también puede configurar tiempos de espera \cite{grpcdeadlines}. El plazo de 1,5 s es una decisión de diseño del equipo.

    \item \textbf{R1 y R6. Notificaciones fuera del flujo crítico.}
    El aviso no es necesario para completar la operación. Si fuera síncrono, una caída del proveedor de SMS retrasaría o rechazaría transferencias y retiros. Por eso los eventos se publican después de confirmar el resultado correspondiente y notificaciones los consume cuando puede. En transferencias se exige confirmación bancaria; en retiros se comprueba también la entrega de efectivo, y en reversos la devolución confirmada. Se usa publicación/suscripción porque el mismo evento alimenta a notificaciones y a auditoría sin que el publicador los conozca \cite{hohpe2003}. Se eligió RabbitMQ sobre Kafka porque el volumen del enunciado, hasta 300 transacciones por segundo, no exige un registro de eventos distribuido a gran escala; es un criterio del equipo, y Kafka sería la opción para reprocesar historia.

    \item \textbf{R2 y F3. Entrega sin pérdidas y sin duplicados.}
    Los mensajes son persistentes, el publicador espera la confirmación del broker y los consumidores confirman tras procesar \cite{rabbitmqconfirms}. La entrega es al menos una vez, por lo que puede haber duplicados. Cada evento lleva \texttt{idOperacion} e \texttt{idEvento}, y los consumidores ignoran los repetidos \cite{featonbyretries}. Así el socio no recibe dos avisos por una misma operación.

    \item \textbf{R5. Auditoría caída.}
    Los eventos esperan en \texttt{auditoria.q} y se procesan al recuperarse. Si el broker no responde, el servicio publicador guarda el evento en su bandeja de salida junto con el estado y reintenta después \cite{awsoutbox}. La cola y la bandeja tienen espacio finito, por lo que se dimensionan y se alerta por acumulación. Si el registro persistente no está disponible, no se inicia un nuevo movimiento. Si el core ya lo ejecutó y falla la escritura local, se recupera y concilia el resultado; no se informa un rechazo definitivo sin comprobarlo.

    \item \textbf{R6. Disponibilidad.}
    Notificaciones y auditoría quedan fuera de la ruta crítica, así que su caída no detiene transferencias ni retiros. El servicio de cuentas no guarda estado y puede replicarse.
\end{enumerate}

\subsubsection{Alternativas descartadas}

\begin{itemize}
    \item REST con JSON sobre HTTP/1.1 para consultar saldo: se prefirió el contrato RPC y la gestión de plazos de gRPC; REST también admite tiempos de espera configurados en el cliente.
    \item Mensajería para consultar saldo: añade un salto mientras alguien espera.
    \item Notificación síncrona dentro de la operación: una caída del proveedor frenaría las operaciones.
    \item Kafka: el volumen del enunciado no lo justifica.
    \item Saldo guardado en el servicio de cuentas: podría no reflejar un débito reciente.
\end{itemize}

\section{Transferencia interna de Ahorro Futuro}

\subsection{Notas del bloque B}

\subsubsection{Interacciones del diagrama}

\begin{enumerate}
    \item \textbf{B1. Cliente móvil/web compartido (A) al Servicio de transferencias.}
    El cliente envía por REST sobre HTTPS el \texttt{idOperacion}, las cuentas de origen y destino y el monto. Conserva el identificador al repetir el intento. El servicio devuelve el identificador y el estado. Confirma el éxito después de verificar el movimiento en el core.

    \item \textbf{B2. Servicio de transferencias al Registro de operaciones.}
    Antes de llamar a antifraude, registra por SQL sobre TCP con TLS la clave (socio, \texttt{idOperacion}), los datos y el estado. Una restricción de unicidad y una inserción atómica permiten que una sola solicitud obtenga la clave. Si se repiten los datos, se devuelve el estado existente. Si cambian, se rechaza el uso de esa clave \cite{featonbyretries}.

    \item \textbf{B3. Servicio de transferencias a Antifraude (C).}
    RPC síncrona mediante gRPC sobre HTTP/2. Envía el \texttt{idOperacion} y los datos de riesgo. Recibe la decisión y su motivo. Sigue el plazo de 200 ms de C \cite{grpcdeadlines}. Ante un rechazo registra el motivo. Ante un error o vencimiento queda pendiente de antifraude. El movimiento en el core requiere aprobación.

    \item \textbf{B4 y B5. Servicio de transferencias al Adaptador y al Core heredado.}
    Tras la aprobación, el servicio envía el identificador, ambas cuentas y el monto al adaptador por gRPC. El adaptador solicita por la interfaz nativa del core el débito y el crédito en una sola transacción. El core devuelve el resultado y la referencia oficial. Si se pierde la respuesta, se consulta o concilia esa referencia antes de reintentar.

    \item \textbf{B6a. Resultado y evento en el Registro de operaciones.}
    El servicio guarda en una misma transacción SQL el estado y el evento pendiente en la bandeja de salida. Conserva \texttt{idEvento}, \texttt{idOperacion}, estado, fecha y referencia del core cuando exista. Si el core confirma y falla la escritura local, la conciliación recupera el resultado oficial y registra el evento con un identificador estable \cite{awsoutbox}.

    \item \textbf{B6b y B6c. Publicador de outbox y RabbitMQ compartido (A).}
    El publicador lee eventos pendientes por SQL y los envía por AMQP 0-9-1 al exchange operaciones, fuera del flujo principal. Espera confirmación del broker y reintenta conservando el \texttt{idEvento}. Auditoría deduplica por ese identificador y confirma después de procesar. Las notificaciones de éxito requieren que el core haya confirmado el movimiento \cite{awsoutbox}.

\end{enumerate}

\subsubsection{Justificación del diseño}

\begin{enumerate}
    \item \textbf{R2 y F3. Consistencia y doble envío.}
    El registro distingue dos operaciones válidas de un doble clic. El core debe admitir una referencia única consultable para conciliar sin repetir el débito. Si sigue sin conocerse el resultado, se mantiene pendiente y se deriva a conciliación antes de considerar un reintento \cite{featonbyretries}.

    \item \textbf{R3 y F2. Antifraude fuera de servicio.}
    La transferencia queda pendiente de evaluación cuando se agota el plazo de 200 ms. Puede retomarse con el mismo identificador cuando antifraude esté disponible.

    \item \textbf{R5. Auditoría y notificaciones.}
    La bandeja de salida cubre una caída del broker y la cola durable de A conserva los eventos cuando cae auditoría. El envío directo sin guardar el evento dejaría operaciones sin registro después de una interrupción \cite{awsoutbox}.

    \item \textbf{R1 y R6. Rendimiento y disponibilidad.}
    Se mediría si, bajo 300 transacciones por segundo, al menos el 95 \% llega a estado final en menos de 2 segundos. Un estado pendiente queda fuera de las completadas. Se propone replicar el servicio y su registro con conmutación ante fallos, y verificar el 99,95 \% con pruebas de caída y recuperación del core, del registro y de antifraude.

\end{enumerate}

\section{Transferencia interbancaria y antifraude de Ahorro Futuro}

\subsection{Notas del bloque C}

\subsubsection{Interacciones del diagrama}

\begin{enumerate}

    \item Servicio de transferencias a Registro de operaciones.
    Consulta y escritura síncronas mediante un controlador SQL
    sobre TCP con TLS. El servicio registra la clave compuesta
    por socio e \texttt{idOperacion}, los datos de la solicitud
    y el estado de la operación. Una restricción de unicidad y
    una inserción atómica deben impedir que dos solicitudes
    reserven la misma clave. Si se repite una solicitud con
    los mismos datos, se devuelve el estado o resultado
    registrado. Si se reutiliza el identificador con datos
    diferentes, se propone rechazar la solicitud. Esta reserva
    corresponde al identificador de la operación, no a una
    reserva de fondos. El motor SQL y su controlador quedan
    por definir.

    \item Servicio de transferencias a Servicio antifraude.
    RPC síncrona mediante gRPC sobre HTTP/2. El orquestador
    envía \texttt{idOperacion} y los datos necesarios para
    evaluar el riesgo. El servicio antifraude aplica las
    reglas configuradas y produce un dictamen. La elección
    de RPC permite expresar la interacción como una llamada
    remota con parámetros y resultados definidos, aunque
    deben contemplarse los fallos propios de la comunicación
    entre procesos \cite{birrell1984}.
    El plazo de 200 ms es una propuesta del diseño y debe
    validarse mediante pruebas.

    \item Servicio antifraude a Servicio de transferencias.
    Respuesta de la misma llamada RPC con decisión, motivo
    e \texttt{idOperacion}. Si el dictamen es aprobado,
    el orquestador guarda el resultado y continúa con la
    solicitud de reserva de fondos. Si es rechazado, registra
    el motivo y no envía la transferencia al banco. Si ocurre
    un error o vence el plazo, registra la operación como
    pendiente de antifraude y suspende el envío. El vencimiento
    puede ser detectado por el cliente sin recibir una
    respuesta válida. La ausencia de respuesta no se
    interpreta como un rechazo por riesgo.

    \item Servicio de transferencias a Núcleo bancario
    heredado.
    Petición y respuesta síncronas mediante la interfaz nativa
    del core. Después de la aprobación antifraude, el
    orquestador envía los datos necesarios para consultar
    la cuenta y solicitar la reserva de fondos. El core
    devuelve el resultado, \texttt{idOperacion} y
    \texttt{referenciaCore}. Solo se continúa hacia la
    pasarela cuando la reserva ha sido aceptada. Si el core
    rechaza la solicitud o no puede determinarse el resultado,
    se suspende el envío. El protocolo y el middleware deben
    precisarse según la interfaz disponible.

    \item Servicio de transferencias a Pasarela interbancaria.
    Comunicación síncrona mediante sockets TCP con TLS y
    una API de sockets nativa. El servicio envía una trama
    con longitud, \texttt{idOperacion}, cuentas de origen y
    destino y monto. El envío requiere aprobación antifraude
    y reserva de fondos aceptada. El contrato de comunicación
    debe definir el formato de los campos, los límites de
    longitud y los códigos de respuesta.

    \item Pasarela interbancaria a Red interbancaria o banco
    destinatario.
    Solicitud y respuesta síncronas mediante sockets TCP
    con TLS. La pasarela adapta la solicitud al protocolo
    acordado con el sistema externo y recibe un acuse de
    recibo, una aceptación o un rechazo, según el contrato.
    La respuesta incluye una referencia cuando corresponda.
    Debe distinguirse entre recepción técnica, aceptación
    para procesamiento y confirmación definitiva. Un acuse
    de recepción no se considera suficiente para registrar
    éxito definitivo, salvo que el contrato establezca
    expresamente ese significado.

    \item Pasarela interbancaria a Servicio de transferencias.
    Retorno del resultado sobre la conexión de la
    interacción 5. El orquestador registra la operación
    como confirmada, rechazada por el banco o por conciliar.
    Si se pierde la respuesta después del envío, no se
    presupone que la operación fracasó. Se consulta o
    concilia el estado mediante las referencias disponibles
    antes de reintentar. Una llamada remota puede dejar
    incertidumbre sobre la ejecución cuando falla la
    comunicación, por lo que el tratamiento de fallos
    forma parte del diseño \cite{birrell1984}.

    \item Servicio de transferencias a Registro de operaciones.
    Escritura síncrona mediante SQL sobre TCP con TLS.
    El servicio guarda el resultado y el evento pendiente
    en un outbox dentro de la misma transacción local.
    Se conservan \texttt{idOperacion}, estado, motivo e
    \texttt{idEvento}. El objetivo es dejar registrado el
    trabajo de publicación junto con el cambio de estado.
    Esta transacción no incluye automáticamente al core
    ni al banco destinatario. La separación entre
    transacciones locales y comunicación entre sistemas
    requiere conservar estado y coordinar el procesamiento
    posterior de los mensajes \cite{helland2017}.

    \item Publicador de outbox a Registro de operaciones
    y RabbitMQ.
    El publicador consulta los eventos pendientes mediante
    acceso SQL síncrono y los publica mediante AMQP y
    RabbitMQ. La publicación es asíncrona respecto del
    flujo principal de la transferencia. Se propone utilizar
    mensajes persistentes, colas durables y confirmaciones
    del publicador para seguir la aceptación de los mensajes
    por el broker. Si no se obtiene confirmación, se gestiona
    el reintento sin asumir que la primera publicación
    necesariamente fracasó. El tratamiento de mensajes
    repetidos forma parte de la coordinación entre sistemas
    independientes \cite{helland2017}.

    \item RabbitMQ a Auditoría y notificaciones.
    Entrega asíncrona mediante AMQP. El consumidor deduplica
    por \texttt{idEvento}, procesa el resultado y envía
    el ACK después de completar el procesamiento previsto.
    El registro de deduplicación debe coordinarse con los
    efectos del consumidor para evitar que una repetición
    produzca efectos indebidos. Este diseño no presupone
    una garantía automática de procesamiento exactamente
    una vez \cite{helland2017}.
    Las notificaciones de éxito definitivo se generan
    únicamente cuando el resultado bancario permite
    establecer dicho estado.

\end{enumerate}

\subsubsection{Justificación del diseño}

\begin{enumerate}

    \item Latencia y respuesta inmediata.
    Antifraude, reserva de fondos e intercambio con la
    pasarela se resuelven de forma síncrona porque sus
    resultados condicionan la continuación del flujo.
    El plazo de 200 ms limita la espera propuesta para
    antifraude, pero no representa la duración total de
    la transferencia. La latencia completa debe medirse
    incluyendo persistencia, evaluación de riesgo, core,
    pasarela y respuesta bancaria.

    \item Elección de RPC y sockets.
    RPC se utiliza para representar la evaluación antifraude
    como una operación remota con parámetros y resultados
    definidos \cite{birrell1984}. gRPC sobre HTTP/2 es la
    tecnología propuesta para implementarla. Los sockets
    TCP con TLS se reservan para el intercambio explícito
    de tramas con la pasarela y la red interbancaria.
    La elección permite implementar ambos mecanismos,
    aunque no demuestra que uno tenga mejor rendimiento.
    Su comparación requiere pruebas bajo condiciones
    equivalentes.

    \item Consistencia e identificación de solicitudes.
    El registro persistente conserva la relación entre
    el identificador, los datos de la solicitud y su estado.
    La restricción de unicidad evita reservar dos veces
    la misma clave, pero no basta por sí sola para impedir
    ejecuciones externas duplicadas. El orquestador también
    debe controlar las transiciones de estado y utilizar
    mecanismos de identificación o consulta admitidos
    por el core y la red. La comunicación entre entidades
    con transacciones independientes requiere tratar
    explícitamente las repeticiones y la evolución
    del estado \cite{helland2017}.

    \item Coordinación con el core.
    La reserva de fondos y el envío interbancario no se
    consideran una única transacción SQL. Se propone
    coordinar sus resultados y definir las acciones
    posteriores: confirmar el movimiento, liberar la
    reserva o mantener la operación pendiente de
    conciliación. El modelo de sagas proporciona un
    fundamento para organizar secuencias de transacciones
    y acciones compensatorias cuando no se utiliza una
    sola transacción global \cite{garciamolina1987}.
    Esto no significa que una transferencia ya confirmada
    por el banco pueda deshacerse mediante una simple
    reversión local.

    \item Tolerancia a fallos.
    El outbox conserva eventos pendientes para permitir
    su recuperación después de una interrupción. Los
    reintentos deben acompañarse de identificación de
    mensajes, deduplicación y control de estado.
    Esta separación permite coordinar transacciones
    locales y mensajes sin suponer que todos los
    participantes comparten una operación atómica
    \cite{helland2017}.

    \item Volumen de mensajes.
    Auditoría y notificaciones quedan fuera de la espera
    principal de la transferencia. Deben medirse la
    capacidad del publicador, el ritmo de consumo y la
    acumulación de eventos. El diagrama no demuestra una
    capacidad concreta de transacciones por segundo.
    Cualquier objetivo de rendimiento debe verificarse
    mediante pruebas de carga y recuperación después
    de fallos.

    \item Alternativas descartadas.
    No se continúa hacia la pasarela sin aprobación
    antifraude ni reserva aceptada. No se confirma el éxito
    a partir de un acuse que solo indique recepción.
    Tampoco se reintenta una operación de resultado incierto
    sin consultar o conciliar su estado. Para auditoría se
    propone publicación con registro persistente en lugar
    de depender únicamente de un envío sin recuperación.
    La persistencia del estado y el tratamiento de mensajes
    repetidos son necesarios cuando la comunicación conecta
    transacciones independientes \cite{helland2017}.

\end{enumerate}



\section{Retiro en cajero automático de Ahorro Futuro}

\subsection{Notas del bloque D}

El bloque D coordina el débito bancario, la entrega física de efectivo y su registro. No se entrega efectivo si la autorización válida no llega dentro de 15 segundos. El reverso requiere distinguir una entrega fallida de una confirmación perdida. Las medidas adicionales se presentan como propuestas de diseño que deben validarse con el equipo.

\subsubsection{Interacciones del diagrama}

\begin{enumerate}
    \item \textbf{D1. Cajero automático al Servidor de cajeros.}
    El cajero envía por socket TCP/TLS el mensaje fijo con operación, tarjeta, monto y código de cajero. Se propone incorporar \texttt{idOperacion} como correlación y conservarlo en los reintentos; no reemplaza los campos exigidos. Repetir una solicitud no debe producir otro débito \cite{featonbyretries}. El formato, tamaño y codificación de los campos deben acordarse.

    \item \textbf{D2. Servidor de cajeros al Cajero automático.}
    Devuelve autorización o rechazo asociado a \texttt{idOperacion}. Solo autoriza después de confirmar el débito en el core. El cajero valida la operación y el plazo antes de iniciar la dispensación; una autorización tardía no habilita la entrega.

    \item \textbf{D3. Cajero automático al Servidor de cajeros.}
    Informa el resultado del dispensador y la evidencia local del retiro. Se propone registrar de forma persistente el intento y su resultado para recuperar la información tras una desconexión o reinicio. El envío del mensaje no sustituye la comprobación de la entrega física.

    \item \textbf{D4. Servidor de cajeros al Adaptador y al Core heredado.}
    Solicita el débito con \texttt{idOperacion} y recibe el resultado y la referencia del core por su interfaz nativa. El plazo previsto es de 8 s, dentro de los 15 s del cajero. Si vence, el débito queda de resultado incierto hasta consultarlo o conciliarlo; el vencimiento no demuestra que no se ejecutó. El protocolo servidor--adaptador debe precisarse con el equipo.

    \item \textbf{D5. Servidor de cajeros al Core para el reverso.}
    Si existe débito y se verifica que no hubo entrega, el servidor solicita su reverso mediante el adaptador. Se propone distinguir la acción débito de la acción reverso dentro de \texttt{idOperacion}: repetir el reverso no debe generar créditos adicionales \cite{featonbyretries}. Solo se marca el reverso como confirmado después de comprobar el resultado del core.

    \item \textbf{D6. Servidor de cajeros al Registro de retiros.}
    Guarda el estado y el evento pendiente en una misma transacción local de la bandeja de salida \cite{awsoutbox}. Se propone conservar \texttt{idOperacion}, \texttt{idEvento}, cajero, monto, fecha, \texttt{referenciaCore} y evidencia de entrega. Esta transacción no incluye automáticamente al core ni al dispensador. El estado debe registrarse durante las transiciones necesarias, no únicamente al terminar el retiro.

    \item \textbf{D7. Publicador de outbox y Registro de retiros.}
    El publicador recupera eventos pendientes y conserva sus identificadores durante los reintentos. La flecha del diagrama representa el paso de datos hacia el publicador, que inicia la consulta. Deben recuperarse los pendientes después de reiniciar \cite{awsoutbox}.

    \item \textbf{D8. Publicador de outbox al RabbitMQ compartido (A).}
    Publicación asíncrona por AMQP 0-9-1 al exchange \texttt{operaciones}. Se proponen mensajes persistentes, colas durables y confirmaciones del publicador \cite{rabbitmqconfirms,rabbitmqreliability}. La confirmación del broker no significa que auditoría ya procesó el evento. El consumidor de A deduplica por \texttt{idEvento} y confirma después de procesar \cite{rabbitmqconfirms}.
\end{enumerate}

El orden normal es D1, D4, D2, dispensación y D3. D6 registra las transiciones necesarias; D7 y D8 publican fuera de la espera principal. D5 se ejecuta únicamente cuando corresponde el reverso. Los números identifican conexiones y no una secuencia lineal del 1 al 8.

\subsubsection{Justificación del diseño}

\begin{enumerate}
    \item \textbf{R4. Límite de 15 s y comprobación de entrega.}
    Se propone contar el plazo desde el envío de D1 mediante un reloj monotónico. Antes de iniciar la entrega, el cajero comprueba que la autorización corresponda a la operación vigente y no haya vencido. Los 8 s del core dejan margen para la red y el procesamiento, pero no garantizan el plazo total. Se contrasta el registro local, el resultado del dispensador y el estado del core. Una autorización o un ACK de TCP no prueban que el socio recibió efectivo; la evidencia disponible depende del dispositivo.

    \item \textbf{F1. Respuesta perdida o tardía y reverso.}
    Al vencer el plazo, el cajero bloquea una nueva entrega de esa operación y registra el vencimiento. El servidor consulta el débito. Si este se confirma y no hubo entrega, solicita D5; si sigue incierto, conserva la operación pendiente y no repite ciegamente el débito \cite{featonbyretries}. Un reverso no encontrado tampoco descarta que aparezca un débito tardío: la conciliación debe resolver ambos resultados.
    \begin{itemize}
        \item Débito confirmado y efectivo no entregado: solicitar reverso idempotente y comprobar su resultado antes de marcar la operación como revertida.
        \item Débito confirmado y entrega completa: mantener el débito y recuperar o reenviar D3, sin volver a dispensar.
        \item Débito o entrega inciertos: marcar la operación como pendiente de conciliación; no afirmar éxito ni reverso definitivo.
        \item Entrega parcial o resultado físico anómalo: conservar importe y evidencia y aplicar conciliación; no asumir un reverso total automático.
    \end{itemize}

    \item \textbf{Idempotencia y recuperación.}
    Se conserva la relación entre identificador, datos y estado. Se rechaza reutilizar el mismo identificador con datos diferentes. Débito y reverso requieren controles independientes y un resultado persistente; el registro local por sí solo no evita repetir efectos en el core \cite{featonbyretries}. Tras reiniciar, el cajero debe recuperar el estado antes de considerar otra dispensación de la misma solicitud.

    \item \textbf{R5. Auditoría y notificaciones.}
    El estado y el evento se guardan juntos en el outbox y se publican con reintentos \cite{awsoutbox}. Si el core confirmó y falla la escritura local, se recupera el resultado oficial y se registra el evento con un identificador estable. Se distingue retiro entregado, reverso confirmado y resultado pendiente. La notificación de éxito requiere entrega comprobada; un reverso se notifica como concluido únicamente tras confirmarlo.

    \item \textbf{R6. Disponibilidad del 99,95 \%.}
    Se propone acordar con B y C redundancia y recuperación del servidor y del registro, plazos acotados y monitorización. El broker único del diagrama es el servicio compartido definido por A, no un segundo broker exclusivo de D. Su despliegue debe ser coherente con las colas replicadas propuestas por A; si existe una sola instancia física, persiste un punto de fallo que debe evaluarse. La durabilidad y recuperación requieren configuración explícita \cite{rabbitmqreliability,rabbitmqquorum}.
    El objetivo se mide sobre un servicio y período acordados. Para un mes de 30 días, el presupuesto temporal se calcula multiplicando 43\,200 minutos por la fracción de indisponibilidad (1 menos 0,9995): resulta 21,6 minutos. Es un cálculo y no un resultado observado. Deben probarse caída del servidor, respuesta perdida del core, reinicio del cajero y caída del broker; todavía no hay resultados que permitan afirmar el cumplimiento de R6.
\end{enumerate}

\subsubsection{Alternativas descartadas}

\begin{itemize}
    \item Repetir el débito con un identificador nuevo: puede originar otro cargo. Se conserva \texttt{idOperacion} y se consulta o concilia el resultado \cite{featonbyretries}.
    \item Aceptar una autorización fuera de los 15 s: contradice R4 y puede coincidir con el tratamiento del reverso. Una respuesta tardía no habilita una nueva dispensación.
    \item Revertir por cualquier confirmación ausente: el efectivo pudo entregarse. La ausencia activa la verificación; el reverso se ejecuta cuando se establece que corresponde.
    \item Publicar sin persistencia previa: una caída entre registro y envío deja una publicación sin recuperar. El outbox conserva el evento junto con el cambio local \cite{awsoutbox}.
    \item Notificar éxito por la confirmación del broker: esta confirma publicación, no entrega física de efectivo ni procesamiento del consumidor \cite{rabbitmqconfirms}.
\end{itemize}

\subsubsection{Coordinación con los demás bloques}

\begin{enumerate}
    \item \textbf{Con A.} Acordar los eventos de retiro entregado, reverso confirmado y resultado pendiente, conservando su relación con la operación original. Auditoría debe reconstruir el proceso; notificaciones debe distinguir los estados confirmados de los pendientes.
    \item \textbf{Con B y C.} Unificar el identificador de operación, el broker compartido y la justificación de R6, sin duplicar nodos de servicios compartidos.
    \item \textbf{Decisiones pendientes.} Definir tamaño y codificación del mensaje fijo, transporte de la correlación, protocolo servidor--adaptador, consulta admitida por el core y evidencia disponible en el dispensador. La descripción inicial ordena el reverso ante confirmación ausente, pero el diseño propone verificar primero si hubo entrega. Esta precisión debe acordarse con A y revisarse con el docente si el contrato del escenario no la resuelve.
\end{enumerate}

\section*{Declaración de uso de inteligencia artificial}

Para preparar las notas y el diagrama del bloque D se utilizó ChatGPT como apoyo para organizar el contenido y consultar documentación técnica. Las propuestas no representan una implementación ni pruebas ejecutadas. La revisión, validación del diseño e integración final corresponden al equipo. Cada integrante debe completar esta declaración con el uso de IA que corresponda a su propio bloque.

\begin{thebibliography}{99}

\bibitem{birrell1984}
A. D. Birrell y B. J. Nelson,
``Implementing remote procedure calls,''
\textit{ACM Transactions on Computer Systems},
vol. 2, núm. 1, pp. 39--59, 1984.
DOI:
\href{https://doi.org/10.1145/2080.357392}
{10.1145/2080.357392}.

\bibitem{grpcdeadlines}
gRPC Authors,
``Deadlines'',
\textit{gRPC Documentation}.
\url{https://grpc.io/docs/guides/deadlines/}.

\bibitem{rabbitmqconfirms}
RabbitMQ,
``Consumer Acknowledgements and Publisher Confirms'',
\textit{RabbitMQ Documentation}.
\url{https://www.rabbitmq.com/docs/confirms}.

\bibitem{rabbitmqquorum}
RabbitMQ,
``Quorum Queues'',
\textit{RabbitMQ Documentation}.
\url{https://www.rabbitmq.com/docs/quorum-queues}.

\bibitem{rabbitmqdlx}
RabbitMQ,
``Dead Letter Exchanges'',
\textit{RabbitMQ Documentation}.
\url{https://www.rabbitmq.com/docs/dlx}.

\bibitem{hohpe2003}
G. Hohpe y B. Woolf,
\textit{Enterprise Integration Patterns: Designing, Building, and Deploying Messaging Solutions}.
Boston, MA, EE.\,UU.: Addison-Wesley, 2003, ISBN 978-0-321-20068-6.

\bibitem{featonbyretries}
M. Featonby,
``Making retries safe with idempotent APIs'',
\textit{The Amazon Builders' Library}.
\url{https://aws.amazon.com/builders-library/making-retries-safe-with-idempotent-APIs/}.

\bibitem{awsoutbox}
Amazon Web Services,
``Transactional outbox pattern'',
\textit{AWS Prescriptive Guidance}.
\url{https://docs.aws.amazon.com/prescriptive-guidance/latest/cloud-design-patterns/transactional-outbox.html}.

\bibitem{helland2017}
P. Helland,
``Life beyond distributed transactions,''
\textit{Communications of the ACM},
vol. 60, núm. 2, pp. 46--54, 2017.
DOI:
\href{https://doi.org/10.1145/3009826}
{10.1145/3009826}.

\bibitem{garciamolina1987}
H. Garcia-Molina y K. Salem,
``Sagas,''
en \textit{Proceedings of the 1987 ACM SIGMOD
International Conference on Management of Data},
pp. 249--259, 1987.
DOI:
\href{https://doi.org/10.1145/38713.38742}
{10.1145/38713.38742}.

\bibitem{rabbitmqreliability}
RabbitMQ,
``Reliability Guide'',
\textit{RabbitMQ Documentation}.
\url{https://www.rabbitmq.com/docs/reliability}.

\end{thebibliography}

\end{document}
